% GENERATED FILE. Edit canonical lessons, then run npm run generate:books.
% canonical-source-hash: fnv1a64:465983f920e32203
% canonical-lessons: TA-C440-macu, TA-C440-currucculal, TA-C440-tuci, TA-C440-kalivu, TA-C440-nekili

\chapter{Dust and Pollution}
\label{ch:ta-dust-and-pollution}
\hlchaptermodality{440}

\begin{chapteropening}
\textbf{By the end of this chapter:} \emph{I can talk about pollution, the environment, dust, waste and plastic.}

Complete the last lesson of chapter 440: I can talk about pollution, the environment, dust, waste and plastic.
\end{chapteropening}

\section[\texorpdfstring{\ta{மாசு} (mācu)}{மாசு (mācu)}]{\ta{மாசு} (mācu) — pollution; impurity}

\label{lesson:TA-C440-macu}



\begin{quote}
Before the new one: say the Tamil for to sprout, then the Tamil for to sway.
\end{quote}

\subsection*{You'll want to know: \ta{மாசு}}
\textbf{\ta{மாசு}} — \emph{mācu} — \textquotedblleft{}pollution; impurity\textquotedblright{}.

\begin{quote}
\ta{நகரங்களில்} \ta{காற்று} \ta{மாசு} \ta{அதிகமாகிவிட்டது}.
\end{quote}

\begin{quote}
\emph{nakaraṅkaḷil kāṟṟu mācu atikamākiviṭṭatu.}
\end{quote}

\begin{quote}
Air pollution has increased in the cities.
\end{quote}

\ta{காற்று} \ta{மாசு} is air pollution.

\begin{cousinweb}
\textbf{\ta{மாசு}} (\emph{mācu}) is native Dravidian (DEDR 4792), \textquotedblleft{}a spot, a stain, a taint, dust\textquotedblright{}.
\end{cousinweb}

\begin{grammarlens}[title={What you've built}]
One more word for getting around.
\end{grammarlens}

\subsection*{Your turn}
\begin{itemize}
\raggedright
  \item \emph{Say it:} \emph{mācu}
  \item \emph{Say it:} \emph{mācu}, once more
  \item \emph{Say it:} \emph{nakaraṅkaḷil kāṟṟu mācu atikamākiviṭṭatu.} — Air pollution has increased in the cities.
  \item \emph{Recall:} say the Tamil for to sprout, then the Tamil for to sway, then say \emph{mācu} again
\end{itemize}

\begin{tcolorbox}[breakable,colback=teal!4,colframe=teal!35!black,title={Before you move on}]
What does \ta{மாசு} mean? (\textquotedblleft{}Pollution; impurity\textquotedblright{}.) Say it once more.
\end{tcolorbox}
\section[\texorpdfstring{\ta{சுற்றுச்சூழல்} (cuṟṟuccūḻal)}{சுற்றுச்சூழல் (cuṟṟuccūḻal)}]{\ta{சுற்றுச்சூழல்} (cuṟṟuccūḻal) — the environment}

\label{lesson:TA-C440-currucculal}



\begin{quote}
Before the new one: say the Tamil for to sway, then the Tamil for pollution.
\end{quote}

\subsection*{You'll want to know: \ta{சுற்றுச்சூழல்}}
\textbf{\ta{சுற்றுச்சூழல்}} — \emph{cuṟṟuccūḻal} — \textquotedblleft{}the environment\textquotedblright{}.

\begin{quote}
\ta{தொழிற்சாலைகள்} \ta{சுற்றுச்சூழலைக்} \ta{கெடுக்கின்றன}.
\end{quote}

\begin{quote}
\emph{toḻiṟcālaikaḷ cuṟṟuccūḻalaik keṭukkiṉṟaṉa.}
\end{quote}

\begin{quote}
Factories damage the environment.
\end{quote}

\ta{சுற்று}, around, + \ta{சூழல்}, surroundings.

\begin{cousinweb}
\textbf{\ta{சுற்றுச்சூழல்}} (\emph{cuṟṟuccūḻal}) is built from \ta{சுற்று}, \textquotedblleft{}around\textquotedblright{} (DEDR 2715), and \ta{சூழல்}, \textquotedblleft{}surroundings\textquotedblright{} (DEDR 2698).
\end{cousinweb}

\begin{grammarlens}[title={What you've built}]
One more word for getting around.
\end{grammarlens}

\subsection*{Your turn}
\begin{itemize}
\raggedright
  \item \emph{Say it:} \emph{cuṟṟuccūḻal}
  \item \emph{Say it:} \emph{cuṟṟuccūḻal}, once more
  \item \emph{Say it:} \emph{toḻiṟcālaikaḷ cuṟṟuccūḻalaik keṭukkiṉṟaṉa.} — Factories damage the environment.
  \item \emph{Recall:} say the Tamil for to sway, then the Tamil for pollution, then say \emph{cuṟṟuccūḻal} again
\end{itemize}

\begin{tcolorbox}[breakable,colback=teal!4,colframe=teal!35!black,title={Before you move on}]
What does \ta{சுற்றுச்சூழல்} mean? (\textquotedblleft{}The environment\textquotedblright{}.) Say it once more.
\end{tcolorbox}
\section[\texorpdfstring{\ta{தூசி} (tūci)}{தூசி (tūci)}]{\ta{தூசி} (tūci) — dust}

\label{lesson:TA-C440-tuci}



\begin{quote}
Before the new one: say the Tamil for pollution, then the Tamil for the environment.
\end{quote}

\subsection*{You'll want to know: \ta{தூசி}}
\textbf{\ta{தூசி}} — \emph{tūci} — \textquotedblleft{}dust\textquotedblright{}.

\begin{quote}
\ta{சாலையில்} \ta{தூசி} \ta{அதிகமாக} \ta{இருந்தது}.
\end{quote}

\begin{quote}
\emph{cālaiyil tūci atikamāka iruntatu.}
\end{quote}

\begin{quote}
There was a lot of dust on the road.
\end{quote}

\ta{தூசி} \ta{தட்டு} is to dust something off.

\begin{cousinweb}
\textbf{\ta{தூசி}} (\emph{tūci}) is native Dravidian (DEDR 3283), \textquotedblleft{}dust\textquotedblright{}.
\end{cousinweb}

\begin{grammarlens}[title={What you've built}]
One more word for getting around.
\end{grammarlens}

\subsection*{Your turn}
\begin{itemize}
\raggedright
  \item \emph{Say it:} \emph{tūci}
  \item \emph{Say it:} \emph{tūci}, once more
  \item \emph{Say it:} \emph{cālaiyil tūci atikamāka iruntatu.} — There was a lot of dust on the road.
  \item \emph{Recall:} say the Tamil for pollution, then the Tamil for the environment, then say \emph{tūci} again
\end{itemize}

\begin{tcolorbox}[breakable,colback=teal!4,colframe=teal!35!black,title={Before you move on}]
What does \ta{தூசி} mean? (\textquotedblleft{}Dust\textquotedblright{}.) Say it once more.
\end{tcolorbox}
\section[\texorpdfstring{\ta{கழிவு} (kaḻivu)}{கழிவு (kaḻivu)}]{\ta{கழிவு} (kaḻivu) — waste; refuse}

\label{lesson:TA-C440-kalivu}



\begin{quote}
Before the new one: say the Tamil for the environment, then the Tamil for dust.
\end{quote}

\subsection*{You'll want to know: \ta{கழிவு}}
\textbf{\ta{கழிவு}} — \emph{kaḻivu} — \textquotedblleft{}waste; refuse\textquotedblright{}.

\begin{quote}
\ta{தொழிற்சாலைக்} \ta{கழிவு} \ta{ஆற்றில்} \ta{கலக்கிறது}.
\end{quote}

\begin{quote}
\emph{toḻiṟcālaik kaḻivu āṟṟil kalakkiṟatu.}
\end{quote}

\begin{quote}
Factory waste flows into the river.
\end{quote}

From \ta{கழி}, to pass away or discharge, + -\ta{வு}.

\begin{cousinweb}
\textbf{\ta{கழிவு}} (\emph{kaḻivu}) is native Dravidian (DEDR 1356), the noun of \ta{கழி}, \textquotedblleft{}to pass, to be discharged\textquotedblright{}.
\end{cousinweb}

\begin{grammarlens}[title={What you've built}]
One more word for getting around.
\end{grammarlens}

\subsection*{Your turn}
\begin{itemize}
\raggedright
  \item \emph{Say it:} \emph{kaḻivu}
  \item \emph{Say it:} \emph{kaḻivu}, once more
  \item \emph{Say it:} \emph{toḻiṟcālaik kaḻivu āṟṟil kalakkiṟatu.} — Factory waste flows into the river.
  \item \emph{Recall:} say the Tamil for the environment, then the Tamil for dust, then say \emph{kaḻivu} again
\end{itemize}

\begin{tcolorbox}[breakable,colback=teal!4,colframe=teal!35!black,title={Before you move on}]
What does \ta{கழிவு} mean? (\textquotedblleft{}Waste; refuse\textquotedblright{}.) Say it once more.
\end{tcolorbox}
\section[\texorpdfstring{\ta{நெகிழி} (nekiḻi)}{நெகிழி (nekiḻi)}]{\ta{நெகிழி} (nekiḻi) — plastic}

\label{lesson:TA-C440-nekili}



\begin{quote}
Before the new one: say the Tamil for dust, then the Tamil for waste.
\end{quote}

\subsection*{You'll want to know: \ta{நெகிழி}}
\textbf{\ta{நெகிழி}} — \emph{nekiḻi} — \textquotedblleft{}plastic\textquotedblright{}.

\begin{quote}
\ta{கடைகளில்} \ta{நெகிழிப்} \ta{பைகள்} \ta{கொடுக்கக்} \ta{கூடாது}.
\end{quote}

\begin{quote}
\emph{kaṭaikaḷil nekiḻip paikaḷ koṭukkak kūṭātu.}
\end{quote}

\begin{quote}
Shops must not give out plastic bags.
\end{quote}

In speech \ta{பிளாஸ்டிக்}; \ta{நெகிழி} is the coined Tamil word.

\begin{cousinweb}
\textbf{\ta{நெகிழி}} (\emph{nekiḻi}) is a modern coinage built on \ta{நெகிழ்}, \textquotedblleft{}to become loose, to slacken\textquotedblright{} (DEDR 2923): a pliant material.
\end{cousinweb}

\begin{grammarlens}[title={What you've built}]
One more word for getting around.
\end{grammarlens}

\subsection*{Your turn}
\begin{itemize}
\raggedright
  \item \emph{Say it:} \emph{nekiḻi}
  \item \emph{Say it:} \emph{nekiḻi}, once more
  \item \emph{Say it:} \emph{kaṭaikaḷil nekiḻip paikaḷ koṭukkak kūṭātu.} — Shops must not give out plastic bags.
  \item \emph{Recall:} say the Tamil for dust, then the Tamil for waste, then say \emph{nekiḻi} again
\end{itemize}

\begin{tcolorbox}[breakable,colback=teal!4,colframe=teal!35!black,title={Before you move on}]
What does \ta{நெகிழி} mean? (\textquotedblleft{}Plastic\textquotedblright{}.) Say it once more.
\end{tcolorbox}
